\section{系统部署与测试}

\subsection{环境要求}

\begin{table}[H]
  \centering
  \caption{推荐运行环境}
  \label{tab:environment}
  \small
  \begin{tabularx}{0.9\textwidth}{L{3.2cm} Y}
    \toprule
    项目 & 要求或说明 \\
    \midrule
    操作系统 & Windows 10/11，本报告命令以 Windows PowerShell 为例 \\
    Python & 3.10 或以上版本 \\
    基础依赖 & Streamlit、Pandas、NumPy、jieba、scikit-learn、Matplotlib 等 \\
    BERT 依赖 & PyTorch、Transformers、Accelerate \\
    硬件 & 基础功能可使用 CPU；BERT 可在 CPU 上运行，CUDA GPU 主要用于加速 \\
    模型文件 & 传统模型约 1.3 MB；BERT 权重约 711 MB，需要单独携带或下载 \\
    可选服务 & ChatECNU API；未配置时自动使用本地建议模板 \\
    \bottomrule
  \end{tabularx}
\end{table}

Git 仓库保留 BERT 配置、Tokenizer、训练参数和评估指标，但不提交
\mbox{\texttt{model.safetensors}}。若权重不存在，系统会自动降级到 TF-IDF 或规则
后端。页面侧边栏会显示当前实际运行的情感后端，避免把降级结果误认为 BERT 推理结果。

\subsection{安装与启动步骤}

\begin{enumerate}
  \item 在项目根目录创建并激活虚拟环境：

\begin{commandblock}
python -m venv .venv
.venv\Scripts\activate
\end{commandblock}

  \item 安装基础依赖和 BERT 依赖：

\begin{commandblock}
pip install -r requirements.txt
pip install -r requirements-bert.txt
\end{commandblock}

  \item 如需使用大语言模型建议，将 \code{.env.example} 复制为 \code{.env}，填写
        \code{LLM_API_KEY}、\code{LLM_BASE_URL} 和 \code{LLM_MODEL}。不配置 API 也不会影响
        情感分析、课程维度和本地建议生成。

  \item 确认 BERT 模型目录为 \code{outputs/bert_model_final}，或通过
        \code{BERT_MODEL_PATH} 指定其他目录。若没有 BERT 权重，系统自动使用传统模型。

  \item 启动应用：

\begin{commandblock}
streamlit run app.py
\end{commandblock}

  \item 浏览器访问 Streamlit 输出的本地地址，依次检查首页概览、单条分析、批量分析、固定案例
        验证和模型评估页面。
\end{enumerate}

\subsection{源代码组成}

主要内容如下：

\begin{itemize}
  \item \code{app.py}：Streamlit 应用入口及五个功能页面；
  \item \code{src/}：预处理、BERT 与传统模型推理、课程维度、关键词和建议生成模块；
  \item \code{scripts/}：数据构建、消融实验、压力测试和报告图表生成脚本；
  \item \code{tests/} 与 \code{data/test_cases.csv}：自动测试和固定演示案例；
  \item \code{requirements.txt}、\code{requirements-bert.txt} 和
        \code{.env.example}：依赖与配置模板；
  \item \code{models/}：可直接运行的传统模型文件。BERT 权重体积较大，随最终电子材料单独
        提交；缺少该权重时系统会自动回退到传统模型，基本功能仍可运行。
\end{itemize}

\subsection{模型训练与实验复现}

传统模型训练命令如下：

\begin{commandblock}
python -m src.train_model ^
  --data data/processed/bilingual_reviews_train.csv ^
  --model-dir models
\end{commandblock}

BERT 训练命令如下：

\begin{commandblock}
python -m src.train_bert ^
  --data data/processed/bilingual_reviews_train.csv ^
  --model-name bert-base-multilingual-cased ^
  --output-dir outputs/bert_model_final ^
  --metrics-path outputs/bert_metrics_final.json
\end{commandblock}

固定案例消融实验和独立压力测试分别使用：

\begin{commandblock}
python scripts/run_ablation_experiment.py ^
  --cases data/test_cases.csv ^
  --output-dir outputs/reports/final_model ^
  --model-dir models

python scripts/run_stress_test.py ^
  --cases data/stress_test_cases.csv ^
  --output-dir outputs/reports/stress_test_final ^
  --model-dir models
\end{commandblock}

\subsection{测试策略}

项目测试分为三个层次：

\begin{enumerate}
  \item \textbf{单元与集成测试：}覆盖文本清洗、语言识别、分词、情感标准化、长文本窗口、课程
        维度证据、模型优先级、降级路径、LLM 本地兜底、数据脚本和页面数据格式。
  \item \textbf{固定案例：}12 条中英文课程评价，用于验证课堂演示中的情感和课程维度结果。
  \item \textbf{独立压力测试：}48 条复杂表达，用于观察缩写、拼写、否定、讽刺、中英混合和长
        文本下的系统表现。
\end{enumerate}

单元测试运行命令为：

\begin{commandblock}
python -m unittest discover -s tests
\end{commandblock}

\subsection{代表性测试用例与结果}

表~\ref{tab:test-cases} 给出 8 个代表性案例，结果来自当前
完整混合流程的固定案例实验。

{\setlength{\tabcolsep}{4pt}
\begin{longtable}{C{0.8cm} L{6.2cm} C{1.7cm} L{3.2cm} C{1.3cm}}
  \caption{代表性测试用例与结果}
  \label{tab:test-cases}\\
  \toprule
  编号 & 输入评价 & 预期情感 & 预期课程维度 & 结果 \\
  \midrule
  \endfirsthead
  \toprule
  编号 & 输入评价 & 预期情感 & 预期课程维度 & 结果 \\
  \midrule
  \endhead
  1 & 老师讲课很清楚，内容也很有用 &
  积极 & 授课方式、教学内容 & 通过 \\
  2 & 作业太多了，每周都做不完 &
  消极 & 作业任务 & 通过 \\
  3 & 实验环境配置太复杂，经常报错 &
  消极 & 实验实践 & 通过 \\
  4 & 课程内容不错，但是考试范围不太明确 &
  中性 & 教学内容、考试安排 & 通过 \\
  5 & 这门课难度比较大，但学完之后收获很多 &
  中性 & 教学内容、学习收获 & 通过 \\
  6 & 希望老师能多给一些代码示例 &
  中性 & 实验实践、教学内容 & 通过 \\
  7 & The instructor explains concepts clearly and the examples are helpful &
  积极 & 授课方式、学习收获 & 通过 \\
  8 & The assignments are too many and the deadlines are stressful &
  消极 & 作业任务 & 通过 \\
  \bottomrule
\end{longtable}
}

\subsection{结果分析}

固定案例中，TF-IDF 单独通过 7/12；
BERT 单独通过 11/12；
完整流程通过 12/12。
该结果说明课堂演示流程能够稳定跑通。
但由于固定案例包含较多典型规则表达，不能把 12/12 解释为整体准确率。
BERT 改善了英文上下文和中文短句的判断；
规则校正进一步处理了“内容不错，但是考试范围不明确”等混合评价。

48 条独立压力测试中，完整流程通过 34 条。
Accuracy 为 0.7083，Macro-F1 为 0.7217。
中英混合类别通过 6/6；
缩写、拼写、否定和长文本类别均通过 5/6。
这说明针对噪声文本和长文本的设计有效。
讽刺类别仅通过 1/6，是当前最主要的失败来源；
因此报告不以固定案例满分替代独立测试结论。
